\documentclass[10pt]{article}
\usepackage[T1]{fontenc}
\usepackage[utf8]{inputenc}
\usepackage[french]{babel}
\usepackage{amsmath,amssymb,mathrsfs,mathpazo}
\usepackage{geometry}
\geometry{a4paper,margin=2cm}
\usepackage{graphicx,array,tabularx,tabularray,colortbl,multirow,multicol,booktabs,enumitem,pifont,ulem,fancybox,diagbox,cancel}
\usepackage{pstricks,pstricks-add,pst-plot,pst-tree,pst-node,pst-text,pst-func,pst-math,pst-bspline,pst-3dplot,pst-eucl,pst-solides3d}
\let\pscalculate\undefined
\usepackage{pst-bezier}
\usepackage{tikz,pgfplots,tkz-tab}
\usetikzlibrary{arrows,arrows.meta,patterns,calc,positioning,trees,decorations.markings,shapes.geometric,shapes.misc,shapes.arrows,intersections,decorations.pathmorphing,decorations.pathreplacing,backgrounds,shadows,babel}
\pgfplotsset{compat=1.18}
\usepgfplotslibrary{fillbetween}
\usepackage[np]{numprint}
\usepackage{listings,makecell,pas-tableur,twoopt}
\let\barre\undefined
\usepackage[pstricks]{tablvar}
\definecolor{darkgreen}{RGB}{0,100,0}
\providecommand{\red}{\color{red}}
\providecommand{\blue}{\color{blue}}
\providecommand{\green}{\color{green}}
\providecommand{\black}{\color{black}}
\providecommand{\white}{\color{white}}
\providecommand{\gray}{\color{gray}}
\providecommand{\cyan}{\color{cyan}}
\providecommand{\magenta}{\color{magenta}}
\providecommand{\yellow}{\color{yellow}}
\providecommand{\violet}{\color{violet}}
\providecommand{\orange}{\color{orange}}
\providecommand{\euro}{\texteuro}
\providecommand{\EUR}{\texteuro}
\providecommand{\ieme}{\textsuperscript{e}}
\providecommand{\degres}{\ensuremath{^\circ}}
\providecommand{\degree}{\ensuremath{^\circ}}
\providecommand{\vv}[1]{\overrightarrow{#1}}
\providecommand{\Coord}[1]{#1}
\providecommand{\decosix}{}
\providecommand{\starredbullet}{\ensuremath{\bullet}}
\pagestyle{empty}
\setlength{\parindent}{0pt}
\usepackage{adjustbox}\DeclareUnicodeCharacter{2236}{:}\DeclareUnicodeCharacter{2190}{\ensuremath{\leftarrow}}
\providecommand{\R}{\mathbb{R}}
\providecommand{\N}{\mathbb{N}}
\providecommand{\Z}{\mathbb{Z}}
\providecommand{\Q}{\mathbb{Q}}
\providecommand{\D}{\mathbb{D}}
\providecommand{\C}{\mathbb{C}}
\providecommand{\e}{{\mathrm{e}}}
\providecommand{\vect}[1]{\overrightarrow{#1}}
\providecommand{\vectt}[1]{\overrightarrow{\mathrm{#1}}}
\providecommand{\vvt}[1]{\overrightarrow{\mathrm{#1}}}
\providecommand{\np}[2][]{{#2~#1}}
\providecommand{\vv}[1]{{\overrightarrow{#1}}}
\renewcommand{\d}{{\mathrm{d}}}
\providecommand{\degres}{{^\circ}}
\providecommand{\degre}{{^\circ}}
\providecommand{\blue}{\color{blue}}
\providecommand{\red}{\color{red}}
\providecommand{\enskip}{\quad}
\providecommand{\llbracket}{\lbrack\!\lbrack}
\providecommand{\rrbracket}{\rbrack\!\rbrack}
\providecommand{\euro}{{\text{€}}}
\providecommand{\no}{{\text{n°}}}
\providecommand{\interval}[2]{{[#1;#2]}}
\providecommand{\intervFO}[2]{{[#1;#2[}}
\providecommand{\dotfill}{\ldots}

\providecommand{\ds}{\displaystyle}
\providecommand{\textding}[1]{#1}
\providecommand{\Oij}{\(\mathrm{O};\vec{\imath},\vec{\jmath}\)}
\providecommand{\Oijk}{\(\mathrm{O};\vec{\imath},\vec{\jmath},\vec{k}\)}
\providecommand{\Ouv}{\(\mathrm{O};\vec{u},\vec{v}\)}
\providecommand{\intervFF}[2]{\ensuremath{[#1;#2]}}
\renewcommand{\euro}{\ifmmode\text{\texteuro}\else\texteuro\fi}
\DeclareUnicodeCharacter{00B0}{\ifmmode^{\circ}\else\textdegree\fi}
\DeclareUnicodeCharacter{20AC}{\ifmmode\text{\texteuro}\else\texteuro\fi}
% Énoncés : Éducation nationale. Transcription et corrigés : APMEP, Denis Vergès et contributeurs. Découpage et classement : CoopMaths, Sébastien Lozano, Éric Elter et contributeurs. Adaptation de la présentation : Maths Explorer.
% CC BY-SA : https://coopmaths.fr/www/about/
% https://coopmaths.fr/alea/static/e3c/2020/tex/e3c_2020_00_sujet64_4.tex
\begin{document}
\noindent\textit{APMEP / CoopMaths. CC BY-SA. Présentation adaptée par Maths Explorer.}\par\medskip
\medskip

Lors des journées classées \og rouges \fg{} selon Bison Futé, l'autoroute qui relie Paris à Limoges
en passant par Orléans est surchargée.

Lors de ces journées classées \og rouges \fg{}, on a pu observer le comportement des
automobilistes faisant le trajet de Paris à Limoges en passant par Orléans.
\begin{itemize}
\item Pour le trajet de Paris à Orléans, 30\,\% d'entre eux prennent la route nationale, les
autres prennent l'autoroute.

\item Pour le trajet d'Orléans à Limoges :

\begin{itemize}
\item  parmi les automobilistes ayant pris la route nationale entre Paris et Orléans,

40\,\% prennent la route départementale, les autres prennent l'autoroute ;

 \item parmi les automobilistes n'ayant pas pris la route nationale entre Paris et Orléans, 
 
45\,\% prennent la route départementale, les autres prennent l'autoroute.
\end{itemize}
\end{itemize}
On choisit un automobiliste au hasard parmi ceux effectuant, en journée classée rouge, le
trajet Paris - Limoges en passant par Orléans.

On note $N$ l'événement \og l'automobiliste prend la route nationale entre Paris et Orléans \fg{}
et $D$ l'événement \og l'automobiliste prend la route départementale entre Orléans et
Limoges \fg.

Si $A$ est un évènement, on note $\overline{A}$ l'évènement contraire de A.

\medskip

\begin{enumerate}
\item  Recopier sur la copie et compléter l'arbre ci-dessous.
\begin{center}
\psset{nodesepA=0pt,nodesepB=3pt,treesep=0.75,labelsep=0.1pt}
\pstree[treemode=R]{\TR{}}
{\pstree{\TR{$N$~~}\taput{$\dots$}}
	{
	\TR{$D$}\taput{$\dots$}
	\TR{$\overline{D}$}\tbput{$\dots$}
	}
\pstree{\TR{$\overline{N}$~~}\tbput{$\dots$}}
	{\TR{$D$}\taput{$\dots$}
	\TR{$\overline{D}$}\tbput{$\dots$}
	}
}
\end{center}



\item Calculer $p\left(\overline{N}\cap \overline{D}\right)$  et interpréter le résultat.
\item Montrer que la probabilité que l'automobiliste ne choisisse pas la Route Départementale
entre Orléans et Limoges est 0,565.
\end{enumerate}
Lors de ces journées classées \og rouges \fg, on donne les temps de parcours suivants :

Paris - Orléans, par autoroute : 3 heures ;

Paris - Orléans, par nationale : 2 heures ;

Orléans - Limoges, par autoroute : 4 heures ;

Orléans - Limoges, par départementale : 3 heures et demie.
\begin{enumerate}[resume]
\item Recopier et compléter le tableau ci-dessous, qui donne pour chaque trajet, le temps en
heure et la probabilité :

\begin{center}
\begin{tabularx}{\linewidth}{|m{3cm}|*{4}{>{\centering \arraybackslash}X|}} \hline
Évènement 		&$N\cap D$ &$N\cap \overline{D}$ &$\overline{N}\cap D$& $\overline{N}\cap\overline{D}$\\\hline
Temps en heure 	& 5,5 &&&\\\hline
Probabilité 	&0,12&&&\\\hline
\end{tabularx}
\end{center}

\item Calculer l'espérance de la variable aléatoire qui donne la durée du trajet en heure et en
donner une interprétation.
\end{enumerate}

\end{document}
